% !PRINTED 76
\weberpage{76}
%<<CONT
zugleich in $\mathfrak{B}'$ enthalten ist, so kann man $a'$ so wählen, dass $\chi_b(a')$ nicht $= 1$ ist. Dann folgt aber aus (9)
\weq{10}{(\chi_b, r) = 0;}
also der Satz:

\gesperrt{1. Wenn die Periode $\eta$ nicht zu $\mathfrak{A}$, sondern zu einer umfassenderen Gruppe $\mathfrak{A}'$ gehört, dann verschwindet jede Resolvente $(\chi_b, r)$, wenn $b$ eine Zahl aus $\mathfrak{B}$ ist, die nicht in $\mathfrak{B}'$ vorkommt.}

Um nun die Bedingungen für diesen Satz weiter zu verfolgen, müssen wir die Bildungsweise der Charaktere berücksichtigen.

Wir wählen die Bezeichnung so, wie wir sie am Schluss des \S.\,16 eingeführt haben. Dann ist, wenn $\beta_{-1}, \beta_0, \beta_1, \ldots \beta_\mu$ die Indices von $b$ und $\nu_{-1}, \nu_0, \nu_1, \ldots \nu_\mu$ die von $n$ sind,
\weq{11}{\chi_b(n) = \varepsilon_{-1}^{\beta_{-1}\,\nu_{-1}}\, \varepsilon_0^{\beta_0\nu_0}\, \varepsilon_1^{\beta_1\nu_1} \ldots \varepsilon_\mu^{\beta_\mu\nu_\mu}.}

Bezeichnen wir mit $r_0, r_1, \ldots r_\mu$ primitive Einheitswurzeln der Grade $2^\lambda, q_1^{\varkappa_1}, q_2^{\varkappa_2}, \ldots q_\mu^{\varkappa_\mu}$, so können wir jede primitive $m^{\text{te}}$ Einheitswurzel in der Form darstellen (Bd.\,I, \S.\,132)
\weq{12}{r = r_0 r_1 \ldots r_\mu,}
und wenn wir $n_0, n_1, \ldots n_\mu$ aus den Congruenzen bestimmen
\weq{13}{n \equiv n_0 \ (\mathrm{mod}\ 2^\lambda), \quad n \equiv n_1 \ (\mathrm{mod}\ q_1^{\varkappa_1}), \ldots n \equiv n_\mu \ (\mathrm{mod}\ q_\mu^{\varkappa_\mu}),}
so zerfällt $(\chi_b, r)$ nach (5) in das Product der folgenden Summen:
\weq{14}{\overset{n_0}{\Sigma}\, \varepsilon_{-1}^{\beta_{-1}\,\nu_{-1}}\, \varepsilon_0^{\beta_0\nu_0}\, r_0^{n_0}, \quad \overset{n_1}{\Sigma}\, \varepsilon_1^{\beta_1\nu_1}\, r_1^{n_1}, \ldots \overset{n_\mu}{\Sigma}\, \varepsilon_\mu^{\beta_\mu\nu_\mu}\, r_\mu^{n_\mu}.}

Wenn nun $b$ eine Zahl ist, die in $\mathfrak{B}$, aber nicht in $\mathfrak{B}'$ vorkommt, so muss nach dem Satze 1. eine von diesen Summen verschwinden.

Dafür ist aber nach \S.\,17 die nothwendige und hinreichende Bedingung die, dass eine der Congruenzen
\weq{15}{\beta_0 \equiv 0\ (\mathrm{mod}\ 2), \quad \beta_1 \equiv 0\ (\mathrm{mod}\ q_1), \ldots \beta_\mu \equiv 0\ (\mathrm{mod}\ q_\mu)}
befriedigt ist, und zwar eine solche, deren Modul mehrfach in $m$ aufgeht.

Diese Bedingung muss zunächst, wenn die Voraussetzung des Satzes 1. zutrifft, für jede Zahl $b$, die in $\mathfrak{B}$, aber nicht in $\mathfrak{B}'$ vorkommt, erfüllt sein. Ist aber $b'$ eine Zahl in $\mathfrak{B}'$, so ist für jeden ganzzahligen Exponenten $x$ das Product $b'^{x} b$ in $\mathfrak{B}$, aber
%CONT>>
